% ================= 1. EL PROBLEMA =================
\section{El problema: tres fallas estructurales}\label{sec:problema}

El estado actual de la IA jurídica para México puede caracterizarse por tres fallas que se refuerzan mutuamente. Ninguna es conjetural: las tres están medidas en literatura revisada por pares o en evaluaciones propias reproducibles, y las tres apuntan a la misma conclusión arquitectónica ---un sistema legalmente usable para México debe ser \emph{anclado en fuentes mexicanas}, \emph{desplegable localmente} y \emph{alimentado por un corpus que hoy no existe}.

\subsection{Alucinación de derecho: fluida, confiada y errónea}

La literatura sobre alucinación legal es ya robusta. Dahl et al. administraron más de 200,000 consultas verificables a GPT-3.5, GPT-4, Llama 2 y PaLM 2, y encontraron tasas de alucinación entre el 58\% (GPT-4) y el 88\% (Llama 2); al responder sobre el \emph{holding} de una sentencia, los modelos alucinan al menos 75\% del tiempo.\cite{dahl2024} Dos hallazgos del estudio son especialmente relevantes para México: el desempeño \emph{empeora sistemáticamente en jurisdicciones menos prominentes} (tribunales inferiores, materias de baja cobertura), y los modelos carecen de autoconciencia del error ---refuerzan premisas jurídicas incorrectas del usuario en lugar de corregirlas.\cite{dahl2024} El derecho mexicano es, desde la perspectiva de un corpus de entrenamiento dominado por texto anglosajón, la definición exacta de una jurisdicción poco prominente.

La recuperación aumentada mitiga pero no elimina el problema. Magesh et al., en la primera evaluación prerregistrada de herramientas jurídicas comerciales con RAG, midieron tasas de alucinación de 17\% en Lexis+ AI y Ask Practical Law AI, y de 33\% en Westlaw AI-Assisted Research, frente a 43\% de GPT-4; el modo de falla más peligroso que documentan es la respuesta \emph{misgrounded}: una cita a una fuente real que no sustenta la afirmación.\cite{magesh2025,stanfordhai2024} Un ensayo controlado aleatorio con 127 estudiantes de derecho refuerza el punto desde otro ángulo: quienes usaron una herramienta RAG (Vincent AI de vLex) produjeron trabajo con una tasa de alucinación comparable a no usar IA, mientras que quienes usaron un modelo de razonamiento sin RAG produjeron mejor análisis pero introdujeron alucinaciones.\cite{ailawlibrarians2026}

Nuestra propia evaluación reproduce el fenómeno en el contexto mexicano. Sometimos un modelo legal de fundación anglosajón de última generación (referenciado en este documento como Thomson-1.0-Small, véase la Sección~\ref{sec:relacionado}) a una batería de 30 preguntas jurídicas ciudadanas canónicas ---laboral, familiar, consumidor, vivienda---. Respondió 30/30 con fluidez y 93\% de las respuestas citaban legislación. La inspección reveló errores sistemáticos de cita: atribuyó la protección contra despido por embarazo al ``art.~139 LFT'' (el correcto es el 164) y la fórmula indemnizatoria de ``3 meses más 20 días por año'' al art.~48 (es el art.~50). Para un ciudadano que no puede verificar, una cita fluida, confiada y errónea es el peor modo de falla posible: tiene la forma de certeza jurídica sin su contenido.

\begin{alerta}{Modo de falla crítico}
El CPT de dominio enseña la \emph{forma} del razonamiento jurídico ---registro, estructura, referencias institucionales--- pero no garantiza anclaje factual en un sistema jurídico específico. La alucinación de citas \textbf{sobrevive al entrenamiento de dominio}. El anclaje jurisdiccional requiere recuperación sobre fuentes primarias de esa jurisdicción.
\end{alerta}

\subsection{Dependencia de nube frente al secreto profesional}

El abogado mexicano opera bajo un deber jurídico de sigilo con sanción penal: el Código Penal Federal tipifica la revelación de secretos conocidos con motivo del empleo, cargo o puesto (art.~210) y la agrava cuando quien revela es un prestador de servicios profesionales (art.~211: uno a cinco años, multa y suspensión de profesión).\cite{cpf} A ello se suma la nueva Ley Federal de Protección de Datos Personales en Posesión de los Particulares (DOF 20 de marzo de 2025), que impone obligaciones reforzadas sobre datos sensibles ---condición típica de los expedientes jurídicos--- y regula expresamente el tratamiento automatizado de datos con implicaciones directas para sistemas de IA.\cite{lfpdppp,cosio2026}

Subir documentos privilegiados de clientes a una API de LLM extranjera ---que puede retener entradas, revisarlas con humanos o usarlas para entrenamiento, bajo jurisdicciones que además habilitan requerimientos gubernamentales foráneos--- es jurídicamente indefendible para un despacho mexicano. No es una preferencia de despliegue: es una restricción de cumplimiento. Las herramientas dominantes (Harvey, CoCounsel, Lexis+ AI, Vincent) se ofrecen exclusivamente como SaaS en nube.\cite{harvey2026,vlex2018} La apertura de regiones cloud en Querétaro por AWS y Google Cloud\cite{aws2024,gcp2024} resuelve la \emph{residencia} del dato pero no la \emph{soberanía de inferencia}: el modelo sigue siendo una caja negra operada por un tercero, con políticas de retención ajenas. La literatura sobre soberanía de IA distingue cuatro capas ---corpus, entrenamiento, lingüística e inferencia--- y coincide en que la que realmente importa a las empresas reguladas es la última: dónde corre el modelo y quién audita lo que entra y sale.\cite{alia2026} Un modelo soberano desplegable on-premise no es un lujo: es la única arquitectura compatible con el deber de sigilo.

\subsection{El corpus mexicano no existe}

SaulLM-7B se entrenó sobre 30B tokens de derecho anglosajón destilados de 94B tokens crudos;\cite{colombo2024} los modelos industriales de Thomson Reuters o Harvey descansan sobre corpus propietarios angloamericanos de decenas de miles de millones de tokens. Para México no existe equivalente ---no porque los datos no existan, sino porque están dispersos en más de 33 portales judiciales, 32 sistemas estatales de gaceta, una gaceta federal con índices desde 1917,\cite{sidof} y repositorios universitarios, la mayoría detrás de protecciones anti-bot, cadenas de tokens de sesión, bloqueos geográficos y aplicaciones legadas. La literatura de NLP jurídico en español refleja este vacío: los recursos a escala existentes (p.~ej. Legal-ES y la tarea IberLegal) cubren textos administrativos y legislativos \emph{españoles}, no mexicanos.\cite{legales2020}

La consecuencia es un círculo vicioso: sin corpus no hay entrenamiento de dominio; sin entrenamiento de dominio, los modelos alucinan el derecho mexicano; y sin anclaje, la IA jurídica es inutilizable precisamente donde más se necesita ---en un país donde la mitad de la población enfrenta problemas jurídicos y solo 5 de cada 100 llegan a una institución.\cite{wjpmx2019,justiciaabierta2025} Romper ese círculo es el objeto de este proyecto.

\begin{figure}[H]
\centering
\includegraphics[width=0.92\textwidth]{fig_alucinacion.png}
\caption{Tasas de alucinación en consultas jurídicas verificables: modelos generales (58--88\%\cite{dahl2024}) frente a herramientas jurídicas comerciales con RAG (17--33\%\cite{magesh2025}). Ninguna herramienta medida está anclada en fuentes mexicanas.}
\label{fig:alucinacion}
\end{figure}
